% ============================================================
%  MTH636 – Game Theory | Lecture 18
%  Chomp \& Extensive Form Games with Chance
% ============================================================
\documentclass[12pt, a4paper]{article}

% ============================================================
%  MTH636 -- Game Theory
%  Notes Preamble — Color-Coded, TikZ-Enhanced
% ============================================================

% --- Core packages ---
\usepackage[T1]{fontenc}
\usepackage[utf8]{inputenc}
\usepackage{lmodern}
\usepackage[margin=2.5cm, top=2.8cm, bottom=2.8cm]{geometry}
\usepackage{microtype}

% --- Math ---
\usepackage{amsmath, amssymb, amsthm, mathtools}
\usepackage{bm}        % bold math
\usepackage{cancel}    % strike-through in math

% --- Color & Boxes ---
\usepackage[dvipsnames, svgnames, x11names, table]{xcolor}
\usepackage[most]{tcolorbox}
\tcbuselibrary{breakable, skins, theorems}

% --- TikZ & Diagrams ---
\usepackage{tikz}
\usepackage{tikz-qtree}
\usetikzlibrary{arrows.meta, positioning, shapes.geometric, fit, backgrounds,
                decorations.pathreplacing, calc, matrix, trees}
\usepackage{pgfplots}
\pgfplotsset{compat=1.18}

% --- Tables ---
\usepackage{booktabs}
\usepackage{multirow}
\usepackage{array}
\usepackage{tabularx}
\usepackage{makecell}

% --- Typography & Layout ---
\usepackage{enumitem}
\usepackage{parskip}
\usepackage{fancyhdr}
\usepackage{titlesec}
\usepackage{titletoc}
\usepackage{hyperref}
\usepackage{graphicx}
\usepackage{float}
\usepackage{caption}

% --- Color Palette ---
\definecolor{defblue}{HTML}{1A5276}      % Definition header
\definecolor{defbluebg}{HTML}{D6EAF8}    % Definition background
\definecolor{thmgreen}{HTML}{1E8449}     % Theorem header
\definecolor{thmbg}{HTML}{D5F5E3}        % Theorem background
\definecolor{exorange}{HTML}{B7770D}     % Example header
\definecolor{exbg}{HTML}{FEF9E7}         % Example background
\definecolor{warnred}{HTML}{922B21}      % Warning/Key Insight header
\definecolor{warnbg}{HTML}{FADBD8}       % Warning background
\definecolor{noticegray}{HTML}{2C3E50}   % Note header
\definecolor{noticebg}{HTML}{EBF5FB}     % Note background
\definecolor{taketawayyellow}{HTML}{7D6608} % Key Takeaway header
\definecolor{takeawaybg}{HTML}{FDFEFE}   % Takeaway background
\definecolor{accentpurple}{HTML}{6C3483} % Accent / proof
\definecolor{lightgray}{HTML}{F2F3F4}    % shading

% ============================================================
%  TCOLORBOX ENVIRONMENTS
% ============================================================

% -- Definition (Blue) --
\newtcolorbox{defbox}[2][]{
  enhanced, breakable,
  colback=defbluebg, colframe=defblue, coltitle=white,
  fonttitle=\bfseries\sffamily,
  title={Definition: #2},
  attach boxed title to top left={yshift=-2mm, xshift=4mm},
  boxed title style={colback=defblue, rounded corners},
  arc=4pt, boxrule=1.2pt,
  left=6pt, right=6pt, top=8pt, bottom=6pt,
  #1
}

% -- Theorem / Proposition (Green) --
\newtcolorbox{thmbox}[2][]{
  enhanced, breakable,
  colback=thmbg, colframe=thmgreen, coltitle=white,
  fonttitle=\bfseries\sffamily,
  title={Theorem: #2},
  attach boxed title to top left={yshift=-2mm, xshift=4mm},
  boxed title style={colback=thmgreen, rounded corners},
  arc=4pt, boxrule=1.2pt,
  left=6pt, right=6pt, top=8pt, bottom=6pt,
  #1
}

% -- Proposition (Green variant) --
\newtcolorbox{propbox}[2][]{
  enhanced, breakable,
  colback=thmbg, colframe=thmgreen!70!black, coltitle=white,
  fonttitle=\bfseries\sffamily,
  title={Proposition: #2},
  attach boxed title to top left={yshift=-2mm, xshift=4mm},
  boxed title style={colback=thmgreen!70!black, rounded corners},
  arc=4pt, boxrule=1.2pt,
  left=6pt, right=6pt, top=8pt, bottom=6pt,
  #1
}

% -- Example (Orange/Yellow) --
\newtcolorbox{exbox}[2][]{
  enhanced, breakable,
  colback=exbg, colframe=exorange, coltitle=white,
  fonttitle=\bfseries\sffamily,
  title={Example: #2},
  attach boxed title to top left={yshift=-2mm, xshift=4mm},
  boxed title style={colback=exorange, rounded corners},
  arc=4pt, boxrule=1.2pt,
  left=6pt, right=6pt, top=8pt, bottom=6pt,
  #1
}

% -- Key Insight / Warning (Red) --
\newtcolorbox{warnbox}[2][]{
  enhanced, breakable,
  colback=warnbg, colframe=warnred, coltitle=white,
  fonttitle=\bfseries\sffamily,
  title={Key Insight: #2},
  attach boxed title to top left={yshift=-2mm, xshift=4mm},
  boxed title style={colback=warnred, rounded corners},
  arc=4pt, boxrule=1.2pt,
  left=6pt, right=6pt, top=8pt, bottom=6pt,
  #1
}

% -- Note (Gray-Blue) --
\newtcolorbox{notebox}[1][]{
  enhanced, breakable,
  colback=noticebg, colframe=noticegray!60, coltitle=white,
  fonttitle=\bfseries\sffamily,
  title={Note},
  attach boxed title to top left={yshift=-2mm, xshift=4mm},
  boxed title style={colback=noticegray!60, rounded corners},
  arc=4pt, boxrule=0.8pt,
  left=6pt, right=6pt, top=8pt, bottom=6pt,
  #1
}

% -- Key Takeaway --
\newtcolorbox{takebox}[1][]{
  enhanced, breakable,
  colback=takeawaybg, colframe=taketawayyellow, coltitle=white,
  fonttitle=\bfseries\sffamily,
  title={Key Takeaway},
  attach boxed title to top left={yshift=-2mm, xshift=4mm},
  boxed title style={colback=taketawayyellow, rounded corners},
  arc=4pt, boxrule=1.2pt,
  left=6pt, right=6pt, top=8pt, bottom=6pt,
  #1
}

% -- Intuition --
\newtcolorbox{intubox}[1][]{
  enhanced, breakable,
  colback=accentpurple!8, colframe=accentpurple!60, coltitle=white,
  fonttitle=\bfseries\sffamily,
  title={Intuition},
  attach boxed title to top left={yshift=-2mm, xshift=4mm},
  boxed title style={colback=accentpurple!80, rounded corners},
  arc=4pt, boxrule=0.8pt,
  left=6pt, right=6pt, top=8pt, bottom=6pt,
  #1
}

% ============================================================
%  PAYOFF MATRIX HELPER
%  Usage: inside tabular, use \payoff{a}{b} for (a,b)
% ============================================================
\newcommand{\payoff}[2]{$(#1,\;#2)$}
\newcommand{\nash}[2]{\underline{$#1$},\underline{$#2$}}  % underline NE payoffs
\newcommand{\Nset}{\mathbb{N}}
\newcommand{\Rset}{\mathbb{R}}
\newcommand{\Sset}{\mathcal{S}}
\newcommand{\Ui}{u_i}

% ============================================================
%  SECTION FORMATTING
% ============================================================
\titleformat{\section}
  {\large\bfseries\sffamily\color{defblue}}
  {\thesection}{1em}{}[\titlerule]

\titleformat{\subsection}
  {\normalsize\bfseries\sffamily\color{thmgreen!80!black}}
  {\thesubsection}{1em}{}

\titleformat{\subsubsection}
  {\normalsize\itshape\sffamily\color{exorange!80!black}}
  {\thesubsubsection}{1em}{}

% ============================================================
%  HEADER / FOOTER
% ============================================================
\pagestyle{fancy}
\fancyhf{}
\fancyhead[L]{\sffamily\small\color{defblue} MTH636 -- Game Theory}
\fancyhead[R]{\sffamily\small\color{defblue} \leftmark}
\fancyfoot[C]{\sffamily\small\thepage}
\renewcommand{\headrulewidth}{0.4pt}

% ============================================================
%  HYPERREF SETUP
% ============================================================
\hypersetup{
  colorlinks=true,
  linkcolor=defblue,
  urlcolor=thmgreen,
  citecolor=accentpurple,
  bookmarksnumbered=true,
  pdftitle={MTH636 Game Theory Notes},
  pdfauthor={MTH636 Student Notes}
}

% ============================================================
%  MISC HELPERS
% ============================================================
\newcommand{\BR}{\text{BR}}        % Best Response
\newcommand{\NE}{\text{NE}}        % Nash Equilibrium
\newcommand{\SPNE}{\text{SPNE}}    % Subgame Perfect NE
\newcommand{\Dom}{\text{dom}}      % Dominance
\newcommand{\mixs}[1]{\sigma_{#1}} % Mixed strategy
\newcommand{\pures}[1]{s_{#1}}     % Pure strategy


% ============================================================
\begin{document}

% Title
\begin{center}
  {\LARGE\bfseries\sffamily\color{defblue} MTH636 -- Game Theory}\\[6pt]
  {\Large\sffamily Lecture 18: Winning Chomp \& Chance Moves}\\[4pt]
\end{center}
\vspace{0.5cm}
\hrule height 1.5pt
\vspace{0.3cm}

% ============================================================
\section{Winning Strategies in Chomp}

Recall Zermelo's Theorem: In any finite, perfect-information, two-player game with no chance moves, if there are no draws, either Player 1 or Player 2 must have a winning strategy.
We apply this to David Gale's game of \textbf{Chomp}. The game cannot end in a draw, as someone must eventually eat the poisoned square $(1,1)$. Therefore, a winning strategy exists. But who has it?

\subsection{Square Boards ($n \times n$)}
\begin{thmbox}{Chomp on a Square Board}
  For an $n \times n$ board ($n > 1$), Player 1 has a winning strategy.
\end{thmbox}

\begin{proof}
  Player 1's winning strategy is explicit and constructive:
  \begin{enumerate}
      \item \textbf{Opening move:} Player 1 chooses the square $(2,2)$. This removes a large block of squares and leaves a perfectly symmetrical board consisting of one row and one column of identical length meeting at the poisoned square $(1,1)$.
      \item \textbf{Subsequent moves:} Player 1 simply mirrors Player 2's moves. If Player 2 chooses a cell $(i, j)$, Player 1 immediately chooses the symmetric cell $(j, i)$.
  \end{enumerate}
  Because Player 1 always restores symmetry, Player 2 is always faced with a symmetric board and is eventually forced to break the symmetry by eating the only square on the diagonal: $(1,1)$. Thus, Player 1 wins.
\end{proof}

\subsection{Rectangular Boards ($n \times m$)}
\begin{thmbox}{Chomp on a Rectangular Board}
  For any $n \times m$ board (where $n \neq m$), Player 1 \textit{always} has a winning strategy.
\end{thmbox}

\begin{proof}
  We prove this by contradiction using a strategy-stealing argument. Because the game is finite and has no draws, Zermelo's theorem dictates that either Player 1 or Player 2 must have a winning strategy.
  
  Suppose for contradiction that Player 2 has a winning strategy $\bar{s}_2$.
  Consider an experimental opening move by Player 1: Player 1 simply eats the single top-right corner square $(n, m)$.
  Because Player 2 possesses a winning strategy $\bar{s}_2$, there must exist some specific response move $(i, j)$ that Player 2 can make following Player 1's opening move, which guarantees a win for Player 2 from that point forward.
  
  However, notice the geometry of Chomp: eating any square $(i, j)$ automatically eats all squares above and to its right, including the top-right square $(n, m)$. 
  This means Player 1 could have simply skipped the "experimental" move and chosen $(i, j)$ on their very first turn! The resulting board state would be exactly identical to the state created after Player 2's "winning" response.
  
  If Player 1 does this, Player 1 effectively "steals" Player 2's winning strategy. Player 1 can now use $\bar{s}_2$ from that point forward to guarantee a win.
  
  Thus, if Player 2 has a winning strategy, Player 1 can steal it to win, which is a logical contradiction (both players cannot have a winning strategy). Therefore, Player 2 cannot have a winning strategy. By Zermelo's theorem, Player 1 \textbf{must} have a winning strategy.
\end{proof}
\textit{Note:} This is a classic non-constructive proof. We know Player 1 wins, but the proof does not tell us exactly what Player 1's opening move should be (unless the board is square).

% ============================================================
\section{Extensive Form Games with Chance}

We now expand our model to include probabilistic elements (like rolling dice or drawing cards). We represent this by introducing a dummy player called \textbf{Nature} (Player 0), who moves randomly according to predetermined probabilities.

\begin{defbox}{Extensive Form Game with Chance Moves}
  A game in extensive form with perfect information and chance moves is a tuple:
  \[ \Gamma = \langle N, V, E, n_0, (V_i)_{i \in N \cup \{0\}}, (p_n)_{n \in V_0}, (u_i)_{i \in N} \rangle \]
  where:
  \begin{itemize}
      \item $N$ is the set of strategic players.
      \item $(V, E, n_0)$ is the directed rooted tree.
      \item $(V_i)_{i \in N \cup \{0\}}$ is a partition of the non-leaf nodes among the agents \textit{and} Nature ($P_0$).
      \item For every node $n \in V_0$ (a node where Nature plays), $p_n$ is a probability distribution over the outgoing edges $A(n)$ from $n$.
      \item $(u_i)_{i \in N}$ are the utility functions mapping leaf nodes to payoffs.
  \end{itemize}
\end{defbox}

When the game reaches a node $n \in V_0$, no strategic choice is made. Instead, nature randomly selects a branch according to the explicit probabilities $p_n$.

\end{document}
